\documentclass[12pt]{article}
\usepackage[margin=1in]{geometry}
\usepackage[utf8]{inputenc}
\usepackage{hyperref}
\usepackage{enumitem}
\usepackage{xcolor}
\usepackage{fancyhdr}
\usepackage{titlesec}
\usepackage{array}
\usepackage{booktabs}
\usepackage{longtable}
\usepackage{fancybox}
\usepackage{parskip}
\setlength{\parskip}{6pt plus 2pt minus 1pt}

% UofSC Color Definitions
\definecolor{garnet}{RGB}{115,0,10}
\definecolor{black}{RGB}{0,0,0}
\definecolor{rose}{RGB}{204,46,64}
\definecolor{atlantic}{RGB}{70,106,159}
\definecolor{congaree}{RGB}{31,65,77}
\definecolor{horseshoe}{RGB}{101,120,11}
\definecolor{neutral90}{RGB}{54,54,54}
\definecolor{neutral70}{RGB}{92,92,92}

% Page setup
\pagestyle{fancy}
\fancyhf{}
\rhead{The Daily Gamecock \textit{\textbf{Opinion}} Editor Onboarding}
\lfoot{\thepage}
\renewcommand{\headrulewidth}{0.4pt}

% Section formatting
\titleformat{\section}[block]{\LARGE\bfseries\color{garnet}}{}{0pt}{}
\titleformat{\subsection}[block]{\Large\bfseries\color{neutral90}}{}{0pt}{}
\titleformat{\subsubsection}[block]{\large\bfseries}{}{0pt}{}

% Hyperlink setup
\hypersetup{
    colorlinks=true,
    linkcolor=atlantic,
    filecolor=atlantic,
    urlcolor=atlantic,
    pdftitle={Opinion Editor Onboarding Guide -- The Daily Gamecock},
    pdfauthor={J.C. Vaught},
    bookmarks=true,
}

% Tip box
\newenvironment{tipbox}[1][Tip]{%
\par\vspace{6pt}\noindent\fcolorbox{atlantic}{white}{\begin{minipage}{\dimexpr\textwidth-2\fboxsep-2\fboxrule}%
\noindent{\color{white}\colorbox{atlantic}{\textbf{#1}}}\par\vspace{4pt}%
}{%
\end{minipage}}\par\vspace{6pt}%
}

% Warning box
\newenvironment{warnbox}[1][Important]{%
\par\vspace{6pt}\noindent\fcolorbox{garnet}{white}{\begin{minipage}{\dimexpr\textwidth-2\fboxsep-2\fboxrule}%
\noindent{\color{white}\colorbox{garnet}{\textbf{#1}}}\par\vspace{4pt}%
}{%
\end{minipage}}\par\vspace{6pt}%
}

% Example box
\newenvironment{exbox}[1][Example]{%
\par\vspace{6pt}\noindent\fcolorbox{horseshoe}{white}{\begin{minipage}{\dimexpr\textwidth-2\fboxsep-2\fboxrule}%
\noindent{\color{white}\colorbox{horseshoe}{\textbf{#1}}}\par\vspace{4pt}%
}{%
\end{minipage}}\par\vspace{6pt}%
}

\begin{document}

\begin{center}
    {\Huge\bfseries\color{garnet} Opinion Editor Onboarding Guide}\\[12pt]
    {\Large The Daily Gamecock}\\[6pt]
    {\large J.C. Vaught}\\[6pt]
    {\normalsize Last updated: January 2026}
\end{center}

\vspace{12pt}
\noindent\rule{\textwidth}{0.6pt}
\vspace{6pt}

\tableofcontents
\newpage

%% ============================================================
\section{Accounts \& Access}
%% ============================================================

You need two logins:

\begin{longtable}{p{3cm} p{5cm} p{6cm}}
\toprule
\textbf{Service} & \textbf{Purpose} & \textbf{Notes} \\
\midrule
\textbf{CEO3} by SNworks & Article writing, previewing, adding titles/authors, managing publication workflow & Get credentials from outgoing editor or EIC \\
\textbf{Slack} (Daily Gamecock workspace) & All editorial communication, budgeting, handoffs & Sign up with your USC \texttt{@mailbox.sc.edu} email rather than \texttt{@email.sc.edu} \\
\bottomrule
\end{longtable}

No other tools are strictly required. Most editors use \textbf{Google Docs} for drafting and collaboration.

\begin{tipbox}[Slack Setup]
Turn on Slack notifications. Put your photo, name, and phone number in your Slack bio.
\end{tipbox}

%% ============================================================
\section{First Reading}
%% ============================================================

Start with \textbf{DG\_editor\_training} (in this same onboarding folder) --- it covers editorial policy and training specific to your role.

The other documents in the style guides folder (\texttt{DG\_Headlines}, \texttt{Writer\_Resources}, \texttt{DG\_statements}) are general advice useful for all writers and worth reading when you have time, but not urgent.

%% ============================================================
\section{Slack Channels}
%% ============================================================

You should be added to the following channels:

\subsection{Core Channels}

\begin{longtable}{p{5.5cm} p{4cm} p{4.5cm}}
\toprule
\textbf{Channel} & \textbf{Who's In It} & \textbf{Purpose} \\
\midrule
\texttt{tdg-opinion} & All opinion writers, managing editors & Main opinion section chat \\
\texttt{tdg-opinionteam} & Managing editors, EIC, opinion leaders, Creative Director & Opinion leadership coordination \\
\texttt{tdg-opinionmultimedia-20XX} & Artists, photographers, illustrators & Visuals for opinion section (name updates yearly) \\
\bottomrule
\end{longtable}

\subsection{Cross-Section Channels}

\begin{longtable}{p{5.5cm} p{8.5cm}}
\toprule
\textbf{Channel} & \textbf{Purpose} \\
\midrule
\texttt{tdg-artsopinion} & Arts \& Culture leaders + Opinion leaders \\
\texttt{tdg-newsopinion} & News leaders + Opinion leaders \\
\texttt{tdg-sportsopinion} & Sports leaders + Opinion leaders \\
\bottomrule
\end{longtable}

\subsection{Organization-Wide Channels}

\begin{longtable}{p{4.5cm} p{5cm} p{4.5cm}}
\toprule
\textbf{Channel} & \textbf{Who's In It} & \textbf{Purpose} \\
\midrule
\texttt{tdg-fullstaff} & Every TDG member & General announcements \\
\texttt{tdg-seniorstaff} & All section leaders/editors + beat writers (paid staff) & Leadership coordination; contains the \textbf{Budget sheet} in bookmarks \\
\texttt{\#ready-for-editing} & All section leaders & Post when a piece is ready for the next stage \\
\bottomrule
\end{longtable}

%% ============================================================
\section{The Budget Sheet}
%% ============================================================

The Budget is a \textbf{Google Sheet} found in the \textbf{bookmarks} at the top of the \texttt{tdg-seniorstaff} channel.

\subsection{Columns You Use}

\begin{longtable}{p{3cm} p{11cm}}
\toprule
\textbf{Column} & \textbf{Format / Notes} \\
\midrule
\textbf{Story} & \texttt{(Opinion) Column: TITLE HERE} --- always prefix with your section \\
\textbf{Author} & Writer's name \\
\textbf{Status} & Current stage: \textit{In writing}, \textit{Up for section editors}, \textit{Up for copy}, etc. \\
\bottomrule
\end{longtable}

\subsection{Color Coding}

\begin{longtable}{p{3cm} p{11cm}}
\toprule
\textbf{Color} & \textbf{Meaning} \\
\midrule
\textbf{Yellow} & On schedule / standard \\
\textbf{Red} & Pushed, delayed, or cancelled \\
\bottomrule
\end{longtable}

Add articles \textbf{at least 7 days} before the target publication date to give Copy and Design enough lead time.

%% ============================================================
\section{CEO3 (SNworks) --- Article Management}
%% ============================================================

\subsection{Getting Started}

\begin{enumerate}
    \item \textbf{Log in} at \url{https://ceo3.getsnworks.com}
    \begin{itemize}
        \item Username: \texttt{sagckopinion@mailbox.sc.edu}
        \item Password: \texttt{Opinion1908}
        \item Select \textbf{CEO3} if prompted (CEO2 is deprecated).
    \end{itemize}
    \item Click \textbf{New} $\rightarrow$ \textbf{New Article} to create a piece.
    \item To find an existing article, use the \textbf{search bar} and type the article name.
\end{enumerate}

\subsection{Setting Up an Article}

\begin{itemize}
    \item \textbf{Title}: Format as \texttt{Column: headline headline}
    \item \textbf{Slug}: Format as \texttt{column-headline-headline-author\_name}
    \item \textbf{Tag}: Set an \texttt{Opinion} tag.
    \item \textbf{Author}: Set the writer as author.
    \item \textbf{Status}: Set to match the current stage (mirrors the Budget sheet).
    \item \textbf{Abstract}: Click the Abstract field to set the \textbf{subheading} (2-sentence summary).
    \item \textbf{Meta tab}:
    \begin{itemize}
        \item \textbf{Infobox} --- set pull quotes here.
        \item \textbf{Social media summary} --- set here as well (Twitter/Instagram-length).
    \end{itemize}
\end{itemize}

\subsection{Editing in CEO3}

\begin{longtable}{p{4.5cm} p{9.5cm}}
\toprule
\textbf{Action} & \textbf{How} \\
\midrule
Add hyperlink & Select text, then \texttt{Cmd+K} (Mac) or \texttt{Ctrl+K} (Windows) \\
Leave a comment for Copy & Select text, click the \textbf{Comment} button, write your note (e.g., ``Source for this claim is in recording ABC'') \\
Leave a comment for the writer & Write \textbf{bold inline text} within the article body, or add a \textbf{bold paragraph}. Optionally sign with your initials (e.g., \textbf{-JC}). \\
Preview & Click the \textbf{Preview} button at the top \\
\bottomrule
\end{longtable}

\subsection{Saving \& Checking In/Out}

\begin{longtable}{p{4.5cm} p{9.5cm}}
\toprule
\textbf{Action} & \textbf{What It Does} \\
\midrule
Save and Check In & Saves the article and locks it (you're done editing for now) \\
Check Out & Re-opens the article for further editing \\
\bottomrule
\end{longtable}

\begin{warnbox}
Always \textbf{check in} before handing off to Copy.
\end{warnbox}

%% ============================================================
\section{Standard Workflow}
%% ============================================================

\subsection{Idea to Budget (Pre-editing)}

\begin{enumerate}
    \item \textbf{Idea surfaces} --- could come from a writer, from you, or from pitch meetings.
    \item \textbf{Add to the Budget} --- add the piece to the Budget sheet in \texttt{tdg-seniorstaff} at least 7 days before target pub date.
    \item \textbf{Submit a Visuals Form} --- this is a workflow in the Senior Staff Slack channel. Ask the writer if they have any specific visual ideas before submitting.
\end{enumerate}

\subsection{Editing (in CEO3)}

We assume the writer knows how to use CEO and has already uploaded their draft. Edit in this order of priority:

\begin{longtable}{c p{3cm} p{8.5cm}}
\toprule
\textbf{Priority} & \textbf{Check} & \textbf{Details} \\
\midrule
A & Factuality & Verify all sources. Do not accept Wikipedia. If a source is not available online, leave a comment in CEO for the Copy desk to verify behind you. \\
B & AP Style & Ensure general AP Style compliance. Opinion pieces have some flexibility, but the basics should be followed. \\
C & Opinion consistency & Watch for the thesis drifting mid-piece. If it shifts, clarify with the writer. \\
D & Readability (optional) & Would a person actually want to read this, or is it dry? \\
E & Add metadata & Add: subhead (abstract), social media summary (Meta tab), pull quote in the Infobox (if applicable). Sometimes add SEO headline and summary. \\
\bottomrule
\end{longtable}

\subsection{Giving Feedback to Writers}

You will typically give feedback \textbf{twice}:

\begin{enumerate}
    \item \textbf{First pass} --- write inline comments in CEO (bold text, optionally signed with initials) or do full rewrites of problem sections. See \texttt{DG\_editor\_training} for comment formatting.
    \item \textbf{Second pass} --- check in the article and message the writer about remaining changes.
\end{enumerate}

\begin{tipbox}[Philosophy]
Do as much of the work yourself as possible. This gives you more control, speeds up the process, and reduces the burden on writers.
\end{tipbox}

\subsection{Handoff to Copy}

\begin{enumerate}
    \item \textbf{Check in} the piece in CEO (save and exit).
    \item \textbf{Update the Budget} --- change the article's status in the Google Sheet.
    \item \textbf{Notify Copy} --- post in \texttt{\#ready-for-editing}, e.g.:\\
    \textit{Column: Article Title is ready for @copy}
    \item \textbf{Wait} for feedback from Copy and Design.
\end{enumerate}

\subsection{After Copy Review}

\begin{itemize}
    \item Copy will reach out to \textbf{you and the writer} if they have additional concerns.
    \item Fix as many issues as you can yourself so the writer has less to deal with. The split is at your discretion.
    \item Often Copy will not require another review if the piece is clean.
    \item If you need to make \textbf{last-minute changes}, let Copy know ASAP so you can coordinate.
\end{itemize}

%% ============================================================
\section{Edge Cases}
%% ============================================================

\subsection{Guest Columns}

\begin{itemize}
    \item \textbf{Less strict} overall: relaxed AP Style, less pressure on opinion consistency and readability.
    \item \textbf{Factuality still applies fully} --- verify sources the same way.
    \item \textbf{Shorter timeline} for visuals, but you still must submit a Visuals Form. Guest columns use a \textbf{standardized image} (no custom illustration).
    \item Pull quote is always \textbf{optional}.
\end{itemize}

\subsection{Fast-Track Columns}

For breaking or time-sensitive events (e.g., a death, mass shooting):

\begin{itemize}
    \item Talk to the \textbf{Editor-in-Chief} to approve an accelerated timeline.
    \item Expect \textbf{no illustrations} --- only simple images.
    \item Visuals Form may not be needed; coordination will happen by word of mouth.
\end{itemize}

\subsection{Pushed (Delayed) Columns}

If a piece needs to be delayed:

\begin{enumerate}
    \item In the Budget sheet, add the text \textbf{PUSHED} (bold) before the article title.
    \item Change the cell color to \textbf{red}.
    \item Move the entry to the \textbf{bottom} of the Budget sheet.
\end{enumerate}

When rescheduling:

\begin{itemize}
    \item Add the article to its \textbf{new date} in the Budget.
    \item Do \textbf{not} write ``PUSHED'' next to the rescheduled entry.
    \item Use the \textbf{standard yellow} cell coloring.
\end{itemize}

If a piece is \textbf{cancelled}: leave it as PUSHED and simply do not re-add it to the calendar.

%% ============================================================
\section{Publication Schedule}
%% ============================================================

\begin{itemize}
    \item Opinion publishes on \textbf{Wednesday} and \textbf{Friday}.
    \item Minimum \textbf{1 piece per day} on each of those days.
    \item Pieces can be either regular columns or guest columns.
    \item Guest columns are preferred when possible --- they can be received and published within a \textbf{2-day window} versus the \textbf{8--9 day window} for regular columns.
\end{itemize}

%% ============================================================
\section{Recruiting Writers}
%% ============================================================

\begin{itemize}
    \item The Daily Gamecock holds \textbf{pitch meetings} organized centrally. You may only get notice about a day beforehand.
    \item For recruiting, \textbf{prioritize prior guest columnists} over brand-new writers. They already know the process and have demonstrated they can deliver.
\end{itemize}

%% ============================================================
\section{Approval Chain}
%% ============================================================

The full approval pipeline for a piece:

\begin{center}
Writer $\rightarrow$ Assistant Editor $\rightarrow$ Editor $\rightarrow$ Copy $\rightarrow$ Managing Editor $\rightarrow$ EIC $\rightarrow$ Photo \& Design
\end{center}

\subsection{Key People to Know}

\begin{longtable}{p{4cm} p{10cm}}
\toprule
\textbf{Role} & \textbf{Why} \\
\midrule
Your writers & You work with them most directly \\
Copy desk & They review everything after you; maintain a good relationship \\
Managing Editor & Approves pieces after Copy; your direct report in the chain \\
Design lead & Coordinates visuals and illustrations for your pieces \\
Editor-in-Chief & Final editorial authority; go to them for fast-track approvals and escalations \\
\bottomrule
\end{longtable}

%% ============================================================
\section{Common Pitfalls}
%% ============================================================

\textbf{Don't dilute the voice.} The most common mistake new opinion editors make is playing it too safe --- softening a writer's raw emotion until the piece loses its punch. Your job is not to neutralize the opinion. Your job is to \textbf{channel and amplify} it while keeping the piece legally safe.

\begin{exbox}
\textbf{Legally unsafe:} ``Student Government hates children and likes cutting costs more.''

Presents opinion as fact; potentially defamatory.

\textbf{Channeled version:} ``Clearly, Student Government seems to love cutting costs, even at the expense of children.''

Same energy, legally defensible as rhetorical opinion.
\end{exbox}

The goal is to preserve the writer's conviction while ensuring the language is protected opinion rather than stated fact. When in doubt, reframe claims as clearly rhetorical rather than removing them.

%% ============================================================
\section{Legal \& University Review}
%% ============================================================

\begin{itemize}
    \item \textbf{SNworks} is only the platform host --- they have no editorial role.
    \item \textbf{The university} generally does not intervene unless:
    \begin{enumerate}
        \item Someone \textbf{complains} about a published piece, OR
        \item You are \textbf{writing about USC itself} (e.g., ``USC is letting nazis speak on campus''). The university will require you to \textbf{get a statement from them} before publication.
    \end{enumerate}
    \item \textbf{DG legal review}: If a piece is particularly aggressive or legally risky, the \textbf{Editor-in-Chief} may send it to legal. This is the EIC's call, not yours.
\end{itemize}

%% ============================================================
\section{Writer Deadlines}
%% ============================================================

\begin{longtable}{p{5cm} p{9cm}}
\toprule
\textbf{Milestone} & \textbf{Minimum Lead Time Before Publication} \\
\midrule
Idea on the Budget & 7--8 days \\
Draft submitted & 3 days (2 days absolute worst case) \\
\bottomrule
\end{longtable}

If a writer misses the draft deadline, the piece will likely need to be pushed (see Section~7.3).

%% ============================================================
\section{Recurring Meetings}
%% ============================================================

\subsection{Sunday}

\begin{longtable}{p{2.5cm} p{4cm} p{7.5cm}}
\toprule
\textbf{Time} & \textbf{Meeting} & \textbf{Who} \\
\midrule
1:00 PM & Senior Staff & All section leaders/editors, paid staff \\
2:00 PM & Full Staff & Every member of The Daily Gamecock \\
2:30 PM & Opinion Section & Opinion writers and opinion leadership \\
\bottomrule
\end{longtable}

\subsection{Wednesday}

\begin{longtable}{p{2.5cm} p{4cm} p{7.5cm}}
\toprule
\textbf{Time} & \textbf{Meeting} & \textbf{Who} \\
\midrule
5:00 PM & Senior Staff & All section leaders/editors, paid staff \\
\bottomrule
\end{longtable}

%% ============================================================
\section{Headline Writing (Opinion)}
%% ============================================================

Opinion headlines \textbf{argue} --- they do not report. The headline must convey WHY it matters and WHO is responsible. See \texttt{DG\_Headlines} in the style guides folder for the full guide.

\subsection{Seven Principles}

\begin{longtable}{p{3.5cm} p{5cm} p{5.5cm}}
\toprule
\textbf{Rule} & \textbf{Do} & \textbf{Don't} \\
\midrule
No colons & Write fluid, declarative sentences & ``USC parking: students pay the price'' \\
No false urgency & Opinions are timeless; skip temporal words & ``just,'' ``now,'' ``recently'' \\
No explainers & Deliver the verdict immediately & ``Why USC's parking policy fails students'' \\
Opinion as fact & State opinions as objective truths & ``Students feel parking is unfair'' \\
No names & Use roles/actions (unless iconic figure) & ``John Smith's policy fails students'' \\
Active blame & Name the actor & ``Mistakes were made'' (passive) \\
So what? & Predict future consequences & Only describing the present problem \\
\bottomrule
\end{longtable}

\subsection{Approved Headline Structures}

\subsubsection{Analysis Pieces}
\begin{enumerate}
    \item Dead/Over Declaration --- ``The era of X is dead''
    \item Comparative Value --- ``X matters more than Y''
    \item Action/Consequence --- ``A acts as / leads to B''
    \item Inverse/Reframe --- Flip the common narrative
\end{enumerate}

\subsubsection{Activism Pieces}
\begin{enumerate}
    \item ``It Is Time'' Ultimatum --- for total overhauls
    \item Mandate --- use ``must,'' not ``should,'' for powerful entities
    \item ``A, Not B'' Responsibility Shift --- push blame up the ladder
\end{enumerate}

%% ============================================================
\section{Editing Tiers}
%% ============================================================

Choose the appropriate level of intervention based on how much work the draft needs:

\begin{longtable}{p{3cm} p{4cm} p{7cm}}
\toprule
\textbf{Tier} & \textbf{When to Use} & \textbf{What You Do} \\
\midrule
Level 1: Full Rewrite & Draft needs complete reconstruction & Rewrite while preserving the writer's research; ask them which original parts to keep \\
Level 2: Sectional Rewrite & Specific sections are weak & Targeted fixes to problem areas only \\
Level 3: Comment Pass & Draft is mostly solid & Strategic inline comments; writer retains full authorship \\
\bottomrule
\end{longtable}

%% ============================================================
\section{Source Verification \& Statements}
%% ============================================================

\subsection{When You MUST Request a Statement}

\begin{itemize}
    \item Criticism or allegations against a person or institution
    \item Claims of negligence, incompetence, or discrimination
    \item Potentially harmful assertions affecting someone's reputation
\end{itemize}

\subsection{Recording Requirements}

Interview recordings \textbf{must} be sent to \textbf{copydesktdg@gmail.com} with the following:
\begin{itemize}
    \item Interviewee name (confirm spelling)
    \item Title/year
    \item Position/major
\end{itemize}

\subsection{Attribution Rules}

\begin{itemize}
    \item Default is \textbf{on the record}. The Daily Gamecock does not accept ``off the record'' without editor approval.
    \item \textbf{Email interviews are not allowed} --- phone or Zoom required.
    \item Email \textbf{statements} are accepted from universities and government only.
    \item Use ``said'' for attribution. Avoid ``stated'' or ``mentioned.''
\end{itemize}

\subsection{If a Source Doesn't Respond}

\begin{itemize}
    \item Follow up \textbf{once} with a polite reminder before your deadline.
    \item Try alternative channels (phone if email fails).
    \item Write: ``did not respond to a request for comment'' --- \textbf{never} ``refused to comment.''
    \item Typical response window: 48--72 hours. Don't delay indefinitely.
\end{itemize}

\subsection{Source Quality (most to least reliable)}

\begin{enumerate}
    \item Government documents / official records
    \item Peer-reviewed academic research
    \item Established news organizations
    \item Official institutional statements
    \item Expert interviews
    \item Personal interviews with affected individuals
    \item Organizational websites
    \item Verified social media accounts
    \item Unverified social media / personal blogs
\end{enumerate}

%% ============================================================
\section{Article Structure}
%% ============================================================

\subsection{Target Word Count}

\begin{longtable}{p{5cm} p{9cm}}
\toprule
\textbf{Type} & \textbf{Words} \\
\midrule
Standard column (online) & 700--900 \\
Print piece & 1,500--1,600 \\
\bottomrule
\end{longtable}

\subsection{Anatomy of an Opinion Piece}

\subsubsection{Introduction}
\begin{itemize}
    \item \textbf{Lede} (1st paragraph): Strong, arguable statement. State the opinion, introduce the topic, convey significance.
    \item \textbf{Nutgraf} (2nd paragraph): Essential context in 75 words max. Answer who, when, where. Must appear within the first three paragraphs.
\end{itemize}

\subsubsection{Body (in any order)}
\begin{itemize}
    \item \textbf{Diagnosis}: Layer the problem with data, statistics, quotes, anecdotes.
    \item \textbf{Prescription}: A solution (direct proposal), alternative (better path), or rationale (explain the flawed logic).
    \item \textbf{Counter-arguments}: Concede valid opposing points, then refute --- explain why your argument is still stronger.
\end{itemize}

\subsubsection{Closing (choose one or combine)}
\begin{itemize}
    \item \textbf{Conclusion}: Restate main point in new words + summarize key evidence.
    \item \textbf{Call-to-Action}: Tell readers what to do (vote, petition, contact, think differently).
    \item \textbf{Coda}: A final reflective thought or anecdote.
\end{itemize}

\subsection{Subhead Rules}

\begin{itemize}
    \item First subhead after the lede must be \textbf{factual, not tonal}.
    \item Place subheads every \textbf{3--4 paragraphs}.
    \item Keep to \textbf{2--7 words}.
    \item Must reflect the content that follows.
\end{itemize}

%% ============================================================
\section{AP Style Quick Reference}
%% ============================================================

\subsection{Dates}
\begin{itemize}
    \item Abbreviate months with dates: \texttt{Aug. 15, 2024} (comma around year)
    \item Never abbreviate \textbf{March, May, June, July} when standing alone
    \item Ranges: \texttt{Oct. 15-18}
    \item Month + year only: \texttt{August 2024} (no comma)
\end{itemize}

\subsection{Names \& Titles}
\begin{itemize}
    \item First reference: full name, title, credentials
    \item Subsequent references: last name only
    \item Student format: ``Jane Doe, a third-year engineering student, said\ldots''
    \item Never call a student a ``major'' (e.g., not ``engineering major'')
    \item Capitalize formal titles before names, lowercase after
    \item Verify spelling via: (1) university directory, (2) LinkedIn + university bio, (3) ask the source directly
\end{itemize}

\subsection{General}
\begin{itemize}
    \item Paragraphs: \textbf{1--4 sentences max}
    \item Periods go \textbf{inside} quotation marks
    \item Use ``said'' for attribution
\end{itemize}

\subsection{Statements Hierarchy}

\begin{longtable}{p{3cm} p{6.5cm} p{4.5cm}}
\toprule
\textbf{Type} & \textbf{Example} & \textbf{Source Required?} \\
\midrule
Definitive & ``USC is increasing housing prices'' & Yes \\
Qualified & ``It seems USC is increasing housing costs'' & Not necessarily \\
Personal & ``I think USC is increasing housing prices'' & No (but discouraged) \\
\bottomrule
\end{longtable}

\end{document}
